\section{DER-SafeAgent}
\label{sec:architecture}

\begin{figure*}[t]
  \centering
  % Figure 1 — DER-SafeAgent final runtime-assurance architecture
% (tag ijcip_final_safeagent_20260810). The visual claim: the LLM does NOT
% own the action surface; a deterministic fast path exists in parallel.
\resizebox{\textwidth}{!}{%
\begin{tikzpicture}[
  font=\small,
  node distance=5mm and 8mm,
  box/.style={draw, rounded corners=1pt, align=center, inner sep=4pt,
              minimum height=8mm},
  det/.style={box, fill=blue!8},
  llm/.style={box, fill=red!10, dashed, thick},
  shield/.style={box, fill=green!10, thick},
  phys/.style={box, fill=gray!12},
  lbl/.style={font=\scriptsize\itshape},
  arr/.style={-latex, thick},
]
% ---- left spine: deterministic decision path -------------------------
\node[det] (ids)  {IDS / SIEM alerts\\+ DER telemetry};
\node[det, below=of ids] (fv)   {Deterministic FeatureView\\
  \scriptsize z-scores, command packets / auth.\ status if available,\\
  \scriptsize sequence regression, staleness, integrity residual\\
  \scriptsize (observable only; no generator ground truth)};
\node[det, below=of fv] (gate) {\textbf{Incident Evidence Gate}\\
  \scriptsize observable protocol evidence, sustained\\
  \scriptsize (frozen thresholds; no model input)};
\node[det, left=22mm of gate] (noop) {monitor / no\_op\\
  \scriptsize no LLM consulted};
% ---- untrusted slow path ---------------------------------------------
\node[llm, right=16mm of fv] (qlora) {\textbf{UNTRUSTED ADVISORY}\\
  QLoRA Hypothesis Agent\\ \scriptsize class + proposed action + rationale\\
  \scriptsize (Qwen2.5-7B / Llama-3.1-8B, $K{=}1$)};
% ---- shield -----------------------------------------------------------
\node[shield, below=14mm of gate, xshift=20mm, text width=7.4cm] (shield)
 {\textbf{RUNTIME-ASSURANCE SHIELD} (deterministic)\\
  \scriptsize fixed 5-primitive registry \; $\cdot$ \; class-aware projection\\
  \scriptsize irreversible proposal $\Rightarrow$ always HITL (action veto)\\
  \scriptsize \textsc{no\_op} vs.\ evidenced incident $\Rightarrow$ act (inaction veto)\\
  \scriptsize emergency fallback ranked by validated estimator (EH)};
\node[box, right=6mm of shield, fill=orange!15] (hitl)
 {Human\\operator};
% ---- bottom: physical layer ------------------------------------------
\node[phys, below=9mm of shield] (act) {Bounded DER action\\
  \scriptsize no\_op / freeze\_setpoint / throttle\_ramp /\\
  \scriptsize request\_ied\_revalidation / (isolate: HITL only)};
\node[phys, below=5mm of act] (grid) {DERMS / feeder\\
  \scriptsize StubFeeder + OpenDSS power flow};
\node[det, left=8mm of act] (fast) {\textbf{Deterministic}\\\textbf{fast path}\\
  \scriptsize immediate bounded\\ \scriptsize response on detection};
% ---- audit ------------------------------------------------------------
\node[box, right=6mm of act, fill=gray!8] (audit) {Append-only\\audit chain};
% ---- edges ------------------------------------------------------------
\draw[arr] (ids) -- (fv);
\draw[arr] (fv) -- (gate);
\draw[arr] (gate) -- node[lbl, left, pos=0.6]{insufficient evidence} (noop);
\draw[arr] (gate.south) -- node[lbl, right, pos=0.35]{evidence threshold met}
  (shield.north west);
\draw[arr, dashed, red!60!black] (gate.east) -| node[lbl, right, pos=0.75]
  {consult} (qlora.south);
\draw[arr, dashed, red!60!black] (qlora.south west) --
  node[lbl, above, sloped]{proposal (untrusted)} (shield.north east);
\draw[arr] (shield) -- (act);
\draw[arr] (shield) -- node[lbl, above]{escalate} (hitl);
\draw[arr] (act) -- (grid);
\draw[arr] (fast) -- (act);
\draw[arr] (fv.west) -| node[lbl, left, pos=0.85]{parallel} (fast.north);
\draw[arr, gray] (act.east) -- (audit.west);
\end{tikzpicture}}

  \caption{\ours{} final architecture. Solid blue components are
    deterministic and own every decision that can reach the feeder; the
    dashed red component is the untrusted QLoRA advisory path, consulted
    only after the Incident Evidence Gate authorizes response under its
    frozen observable-evidence thresholds, and only able to select within the
    admissible candidate set the shield
    computes. A model-proposed irreversible primitive is never
    auto-executed. The deterministic fast path responds immediately in
    parallel; every executed action lands in an append-only audit chain.}
  \label{fig:architecture}
\end{figure*}

\subsection{Design rule and trust boundary}
\label{sec:method-layering}
\ours{} follows the runtime-assurance construction of
Simplex~\cite{Sha2001Simplex} and action
shielding~\cite{Alshiekh2018Shielding}: an unverified, high-capability
component may propose; a small deterministic, directly testable layer decides
what may execute. Three consequences structure the whole design
(Fig.~\ref{fig:architecture}). First, the LLM never emits executable
commands --- it selects from a fixed five-primitive registry
\{\textsc{no\_op}, \textsc{freeze\_\allowbreak setpoint},
\textsc{throttle\_\allowbreak ramp},
\textsc{request\_\allowbreak ied\_\allowbreak revalidation},
\textsc{isolate\_\allowbreak inverter}\}, so a
wrong output is a selection error, not arbitrary code execution. Second,
every quantity that can \emph{relax} a constraint is computed
deterministically from telemetry, never taken from model output; the
adversarial evaluation (\S\ref{sec:rq1}) shows why this rule is
load-bearing. Third, a deterministic fast path produces a bounded
response immediately, so consulting the model never delays containment.

\paragraph{Notation and invariants}
At tick $t$ the detector observes telemetry/packets $o_t$; a deterministic
FeatureView $\phi(o_{1:t})$ yields evidence state $e_t$ (\S\ref{sec:method-featureview}).
The Evidence Gate is a predicate $g(e_{1:t})\in\{0,1\}$
(\S\ref{sec:method-gate}). The registry is a fixed set
$\mathcal{A}=\{\textsc{no\_op},\textsc{freeze},\textsc{throttle},
\textsc{revalidate},\textsc{isolate}\}$ with irreversible subset
$\mathcal{I}=\{\textsc{isolate}\}$. The shield computes an admissible candidate set
$S(e_t)\subseteq\mathcal{A}$; the untrusted model proposes
$a^M_t\in\mathcal{A}\cup\{\bot\}$. The executed action is
$$
a^{\text{exec}}_t=\Pi\big(g(e_{1:t}),\,S(e_t),\,a^M_t,\,\tau_t\big),
$$
where $\Pi$ is the deterministic projection and $\tau_t$ the deadline state.
The design enforces, and we test, the following invariants (unit tests in
\texttt{test\_\allowbreak runtime\_\allowbreak safe\_\allowbreak gate.py} and
\texttt{test\_\allowbreak safety\_\allowbreak projection.py}):
\begin{enumerate}
\item[(I1)] $a^{\text{exec}}_t\in\mathcal{A}$ (only registry primitives execute);
\item[(I2)] $g(e_{1:t})=0 \Rightarrow a^{\text{exec}}_t=\textsc{no\_op}$ (gate closed $\Rightarrow$ monitor only);
\item[(I3)] $a^M_t\in\mathcal{I} \Rightarrow a^{\text{exec}}_t\notin\mathcal{I}$ without human approval (a model-proposed irreversible action is never autonomously executed);
\item[(I4)] if independent evidence $e_t$ establishes an incident, $a^M_t=\textsc{no\_op}$ cannot force $a^{\text{exec}}_t=\textsc{no\_op}$ (model inaction cannot veto evidenced response);
\item[(I5)] a late or failed model call leaves the deterministic fast-path action in force (the model can never \emph{prevent} the bounded response);
\item[(I6)] every proposal, veto, fallback, escalation, and execution is appended to the audit chain.
\end{enumerate}
I3 and I4 are the two-sided rule that model output may gate neither action
nor inaction; the evidence in $g$ and $S$ is computed deterministically,
so an attacker who controls the model's \emph{input or output} cannot
change it \emph{through the model}. This is independence from the LLM, not
robustness to a compromised network: an attacker able to inject or forge
unauthenticated protocol traffic can still move the gate by manufacturing
observable evidence (\S\ref{sec:bg-trust}), which is why the architecture
escalates irreversible actions to a human rather than trusting any single
evidence source. These invariants are deterministic and directly testable;
we do not claim a formal proof.

``Trustworthy'' in this paper denotes implemented, measured properties
--- bounded action generation, schema compliance, containment within the
tested adversarial suite, human oversight, auditability, consequence
awareness, conservative fallback --- not a formal guarantee. The supplementary
material maps each property to its mechanism, the experiment that evaluates it,
and its remaining limitation.

\subsection{Deterministic FeatureView}
\label{sec:method-featureview}
A deterministic first stage converts the rolling telemetry/event window
into a structured FeatureView. Every feature is computed from quantities
observable on the wire or from the defender's grid model, never from
attacker-set ground truth. Table~\ref{tab:provenance} lists each feature,
its observable source, and the attacker's influence over it. The features
are: per-asset z-scores against a 60-step baseline; voltage-band
excursion; a forged-command-packet count; a duplicate-payload signature;
a persistent-freeze flag; a SCADA staleness signature (consecutive
samples whose reported timestamps fail to advance --- the fingerprint of
replay and denial-of-service, which re-emit captured frames); a protocol
sequence-number regression (a frame carrying a sequence number below the
running high-water mark --- the wire-visible signature of replay); and a
telemetry-integrity residual (a reported power above the asset's rated
capacity, physically impossible and therefore a corrupted reading ---
the observable signature of a high-magnitude false-data-injection
attack). No feature is read from attacker-set ground truth; a provenance
audit (Table~\ref{tab:provenance}) records, for each feature, its
observable source and the attacker's influence over it, and we report the
detectability this observable-only discipline permits --- and the FDI
limit it exposes --- in \S\ref{sec:rq2}\@. Every downstream stage, including the LLM prompt,
consumes this single numeric grounding; severity is a bounded
deterministic score, not a model output.

% Table 3: trust boundary and evidence provenance.
% Page-breaking xltabular so the tall table never overflows into the page footer.
{\footnotesize
\begin{xltabular}{\textwidth}{@{}>{\raggedright\arraybackslash}p{0.19\textwidth}
 >{\raggedright\arraybackslash}p{0.22\textwidth}
 >{\raggedright\arraybackslash}p{0.13\textwidth}
 >{\raggedright\arraybackslash}X@{}}
\caption{Evidence provenance and trust boundary for the runtime FeatureView.
Every runtime feature is computed from wire-observable telemetry/packets or
the defender's grid model; generator ground-truth flags are excluded by design
(last row). Wire-observable is not authenticated: an attacker able to inject
unauthenticated traffic can forge much of this evidence, so it is robust to a
compromised model but not a compromised network. ``Attacker influence'' is the
attacker's ability to forge or suppress the feature.}\label{tab:provenance}\\
\toprule
Feature & Observable source & Attacker influence & Runtime use / failure mode \\
\midrule
\endfirsthead
\multicolumn{4}{@{}l}{\emph{Table~\ref{tab:provenance} (continued)}}\\[2pt]
\toprule
Feature & Observable source & Attacker influence & Runtime use / failure mode \\
\midrule
\endhead
\midrule
\multicolumn{4}{r@{}}{\emph{continued on next page}}\\
\endfoot
\bottomrule
\endlastfoot
Command-packet count & DNP3/IEC-61850 command frame on the wire & Can forge command frames & Opens gate; forged commands indistinguishable from legitimate ones absent authentication (benign-command discrimination is a stated limit) \\
\addlinespace[2pt]
Duplicate-payload signature & Repeated identical frame payloads & Can inject duplicates & Replay/echo signature; benign historian echo is an irreducible ambiguity \\
\addlinespace[2pt]
Staleness (non-advancing timestamp) & Frame reported-time field & Can freeze/replay timestamps & DoS/replay signature; a re-timestamped replay evades it (caught instead by sequence regression) \\
\addlinespace[2pt]
Sequence regression & IEC-61850 sqNum / DNP3 sequence vs.\ running high-water mark & Cannot lower a delivered seq without detection unless it also spoofs the counter end-to-end & Replay signature; observable, not authenticated \\
\addlinespace[2pt]
Integrity residual & Reported power vs.\ rated capacity (defender grid model) & Bounded: a value above capacity is physically impossible & High-magnitude FDI; \emph{sub-capacity FDI is not observable here} (documented miss) \\
\addlinespace[2pt]
Persistent freeze & $\ge$10 consecutive zero-power samples & Can force zero output & DoS signature \\
\addlinespace[2pt]
Per-asset z-score, voltage band & Reported power / bus voltage & Can bias reported power (soft) & Soft evidence; not separable from benign transients, so not used to open the gate alone \\
\midrule
\textbf{Excluded by design:} any generator ground-truth flag & \emph{e.g.\ an injector \texttt{tampered} bit} & \emph{fully attacker-/generator-set; an oracle} & \textbf{Never read by the runtime}; unit tests confirm no such field enters a feature \\
\end{xltabular}
\par}


\subsection{Incident Evidence Gate}
\label{sec:method-gate}
The shield constrains \emph{which} action may execute, but our earlier
evaluation showed it had no mechanism for deciding that \emph{no}
incident exists: on benign stochastic days, cloud transients and
measurement noise pushed the pipeline into unnecessary mitigation on
every episode (\S\ref{sec:rq2}). The Evidence Gate closes this gap
deterministically. It opens only when \emph{sustained protocol-level
evidence} --- itself wire-observable and, absent authentication, forgeable
by a network attacker (\S\ref{sec:bg-threat}): a forged command packet, a
duplicate-payload/sequence
regression (replay), non-advancing timestamps or a persistent zero-freeze
(DoS), or a rating-plausibility integrity residual (high-magnitude FDI),
all observable (Table~\ref{tab:provenance}) --- is sustained for the
calibrated persistence window. A physical-only (z-score/voltage) path
exists but its calibrated persistence threshold effectively disables it:
on the development split, no physical-only setting separated attacks from
benign PV transients without either missing attacks or admitting benign
days, which is itself a finding about DER telemetry --- and, after the
observable-only evidence set, is the reason a \emph{sub-capacity} FDI (whose only trace
is such a physical deviation) is not reliably separable from benign
operation without redundant sensing (\S\ref{sec:rq2}). While the gate is
closed the system monitors; no mitigation executes and \emph{the LLM is
not consulted}, which also removes the model's exposure to benign-day
prompt content. Thresholds are calibrated on a development split and
frozen before the held-out evaluation; the gate reads no model output.

\subsection{Untrusted advisory path}
\label{sec:method-advisory}
When the gate opens, a QLoRA-adapted~\cite{Dettmers2023_QLoRA} open-weights
backbone (Qwen2.5-7B~\cite{Hu2024_Qwen25} or
Llama-3.1-8B~\cite{Dubey2024_Llama31}, NF4 4-bit on a single RTX~3080) receives the compact
FeatureView prompt and returns a JSON hypothesis: attack class, affected
asset, confidence, rationale, proposed action. Serving is strict: a run
aborts if the intended backbone or adapter is not loaded, and every call
records the adapter SHA-256, prompt hash, raw output, parse status and
latency, so no deterministic fallback can masquerade as a model result:
a silent fallback would invalidate any claim about model behaviour, so it
is made impossible by construction rather than assumed absent.
At each eligible gated decision tick the advisory path draws one greedy
generation ($K{=}1$); a sustained incident may therefore trigger repeated
consultations across ticks (the per-configuration call counts in
Table~\ref{tab:holdout-v3}). Here $K$ is the number of generations \emph{per
consultation} for self-consistency, not the number of consultations per
incident: $K{=}1$ draws one generation, $K{=}3$ draws three. The default is
$K{=}1$; $K{=}3$ was removed after measurement showed it tripled latency without
a decision-quality gain that survives the shield. Fine-tuning uses a 200-example
DER/ICS corpus with assistant-turn loss masking (supplementary material).

\subsection{Runtime-assurance shield}
\label{sec:method-shield}
The shield computes, deterministically, the \emph{admissible candidate set}
for the current situation and lets the model choose only within it.
Removals are recorded per decision: actions inappropriate for the
evidence class (a telemetry-layer compromise must not be answered by
physically disconnecting a healthy asset), actions ineffective against
the mechanism, and all non-\textsc{no\_op} actions when severity is
below threshold. Two rules are absolute. (i)~A model-proposed
\emph{irreversible} primitive (\textsc{isolate\_inverter}) is never
auto-executed: where the evidence class admits it as a candidate at all it
escalates to a human operator, and where the class forbids it the shield
vetoes it and falls back deterministically
(Table~\ref{tab:action-policy}). An earlier
variant gated this on model-stated confidence and was falsified by
adversarial evaluation --- injected prompts manufacture exactly the
confidence the gate keyed on, making that variant worse than no shield
(\S\ref{sec:rq1}). Attacker-influenceable quantities do not
relax constraints. (ii)~When the model's proposal is unusable ---
invalid, out of registry, below the confidence floor, or vetoed --- the
shield falls back deterministically, ranking the admissible candidate set with the
horizon-aware impact estimator described next.
Table~\ref{tab:action-policy} states the complete frozen policy per evidence
class --- which actions may execute autonomously, which are escalation-only,
which are removed and why, and the deterministic fallback order --- together
with the testbed action semantics and recovery condition. The policy is
author-specified: the evaluation establishes conformance to this stated
policy, not an independent validation of its operational correctness or
completeness.

\begin{table}[t]
\centering\small
\caption{Frozen class-aware action policy of the shield (source of truth:
\texttt{safety\_projection.py}; machine-readable export
\texttt{code/configs/action\_policy.yaml}, verified by
\texttt{tests/test\_action\_policy.py}). Actions: N = \textsc{no\_op},
F = \textsc{freeze\_setpoint}, T = \textsc{throttle\_ramp}, R =
\textsc{request\_ied\_revalidation}, I = \textsc{isolate\_inverter}
(irreversible). \emph{Autonomous} actions may execute without a human; a
model-proposed I is never auto-executed --- escalated to a human where the
class admits it (I3), vetoed with deterministic fallback where it does not
--- and the deterministic fallback never selects it. Non-N actions require
severity $\geq 0.30$; honouring a model choice requires stated confidence
$\geq 0.40$. Removals: inappr.\ = class-inappropriate, ineff.\ =
mechanism-ineffective. Testbed semantics: F holds the last validated
setpoint (70\% of rated kW) for the episode remainder; T restricts ramping
to the IEEE\,1547 minimum; I is absorbing and requires operator
restoration. No automatic release is modelled: stand-down requires the
gate's sustained evidence to clear and an operator/DERMS resume, never
model output alone (I4). The policy is author-specified, motivated by the
class mechanism and counterfactual rollouts; the evaluation establishes
conformance to it, not its independent operational validation.}
\label{tab:action-policy}
\footnotesize\setlength{\tabcolsep}{4pt}
\begin{tabular}{lllll}
\toprule
Evidence class & Autonomous & Escalation-only & Removed & Fallback \\
\midrule
Command spoof & N, F, T & I & R\,(ineff.) & F$\rightarrow$T$\rightarrow$N \\
Replay & N, F, T, R & --- & I\,(inappr.) & R$\rightarrow$F$\rightarrow$N \\
DoS & N, F & I & T,R\,(ineff.) & F$\rightarrow$N \\
FDI (rating residual) & N, F, R & --- & T\,(ineff.); I\,(inappr.) & F$\rightarrow$R$\rightarrow$N \\
No evidence / benign & N & --- & F,T,R,I\,(inappr.) & N \\
\bottomrule
\end{tabular}
\end{table}


\subsection{Horizon-aware impact estimator (EH)}
\label{sec:method-estimator}
The original estimator integrated per-action cost over a fixed 60\,s
window. Counterfactual validation --- branching the simulator at the
decision instant and executing every candidate under identical exogenous
trajectories --- showed it selected a tie-optimal action in only a third of
configurations (Table~\ref{tab:eh-v3}): over the true incident horizon,
candidates differ by two orders of magnitude more than over 60 seconds, so a
60\,s ranking is usually wrong. The repaired estimator EH uses only
decision-time information: a deterministic class-family hypothesis from
protocol evidence (spoof-like, FDI-like, or stale, the latter because
replay and DoS are indistinguishable at decision time), a
remaining-horizon estimate from a development-calibrated mean-residual-life prior,
and restoration-horizon accounting for absorbing actions such as
isolation. The true attack end time is never an input. EH is retained in
the architecture only because it passed a hash-frozen held-out
holdout (\S\ref{sec:rq3}); its point predictions remain
imprecise, so it is used to \emph{rank} the admissible candidate set, never to claim
impact magnitudes.

\subsection{Fast path, slow path, oversight, audit}
\label{sec:method-deployment}
The deterministic fast path (FeatureView $\rightarrow$ evidence class
$\rightarrow$ bounded primitive) responds within one control tick. The
model path is a slow path under an explicit service-level objective:
measured consultation latency (${\approx}5.5$\,s at $K{=}1$ on the reference GPU)
means the model contributes within a 15\,s SLO and is skipped ---
without loss of containment --- under tighter deadlines
(\S\ref{sec:rq4}). Escalations carry the model's proposal,
rationale and the shield's veto reasons to the operator. Every decision
appends to a SHA-256 hash-chained audit record. Relative to a trusted stored
chain head, verification detects modification, deletion, and reordering of
records; it does not provide origin authentication, trusted timestamping, or
protection against complete recomputation by an attacker who controls the entire
log. Full prompts, schemas,
state-machine contracts, and protocol examples are in the supplementary
material.
